\documentclass[12pt,a4paper]{report}
\usepackage[utf8]{inputenc}
\usepackage[french]{babel}
\usepackage[T1]{fontenc}
\usepackage{geometry}
\geometry{hmargin=2.2cm,vmargin=2.5cm}
\usepackage{hyperref}
\usepackage{graphicx}
\usepackage{booktabs}
\usepackage{tabularx}
\usepackage{array}
\usepackage{amsmath,amssymb}
\usepackage{listings}
\usepackage{xcolor}
\usepackage{float}
\usepackage{caption}
\usepackage{fancyhdr}
\usepackage{tikz}
\usetikzlibrary{shapes.geometric, arrows, positioning}

% Couleurs
\definecolor{navy}{RGB}{15, 23, 42}
\definecolor{primaryblue}{RGB}{37, 99, 235}
\definecolor{darkslate}{RGB}{30, 41, 59}
\definecolor{lightgrey}{RGB}{248, 250, 252}
\definecolor{bordergrey}{RGB}{226, 232, 240}
\definecolor{codegreen}{RGB}{22, 163, 74}
\definecolor{codered}{RGB}{220, 38, 38}

\hypersetup{
    colorlinks=true,
    linkcolor=primaryblue,
    citecolor=primaryblue,
    urlcolor=primaryblue,
    pdftitle={Rapport Technique TEK-UP - SecOps K8s AWS},
    pdfauthor={Khayem ZRIBI}
}

\lstdefinestyle{customlatex}{
    backgroundcolor=\color{lightgrey},
    basicstyle=\ttfamily\footnotesize\color{darkslate},
    breaklines=true,
    captionpos=b,
    commentstyle=\color{codegreen},
    frame=single,
    rulecolor=\color{bordergrey},
    keywordstyle=\color{primaryblue}\bfseries,
    stringstyle=\color{codered},
    showstringspaces=false,
    tabsize=2
}
\lstset{style=customlatex}

\pagestyle{fancy}
\fancyhf{}
\fancyhead[L]{\small\textcolor{gray}{TEK-UP | Projet Ingénieur Cybersécurité (2025–2026)}}
\fancyhead[R]{\small\textcolor{gray}{SecOps / Incident Response K8s AWS}}
\fancyfoot[L]{\small\textcolor{gray}{Khayem ZRIBI — Cycle Ingénieur}}
\fancyfoot[R]{\small\thepage}
\renewcommand{\headrulewidth}{0.4pt}
\renewcommand{\footrulewidth}{0.4pt}

\begin{document}

% ==============================================================================
% PAGE DE GARDE
% ==============================================================================
\begin{titlepage}
    \centering
    \vspace*{0.5cm}
    {\Large\textbf{TEK-UP University}}\\[0.2cm]
    {\large\textit{École Supérieure Privée des Technologies de l'Information et du Management}}\\[0.1cm]
    {\small Année Universitaire 2025–2026}\\[1.5cm]

    \rule{\linewidth}{1.5pt}\\[0.4cm]
    {\Huge \bfseries Conception et Déploiement d'une Plateforme Cloud Sécurisée Kubernetes sur AWS avec Détection et Réponse Automatique aux Incidents}\\[0.2cm]
    {\Large \textcolor{primaryblue}{\textbf{SecOps / Incident Response (PICERL) et Machine Learning}}}\\[0.3cm]
    \rule{\linewidth}{1.5pt}\\[1.5cm]

    \begin{minipage}{0.48\textwidth}
        \begin{flushleft} \large
            \textbf{Rédigé par :}\\
            \textbf{Khayem ZRIBI}\\
            2\textsuperscript{ème} Année Cycle Ingénieur\\
            Spécialité : Cybersécurité
        \end{flushleft}
    \end{minipage}
    \hfill
    \begin{minipage}{0.48\textwidth}
        \begin{flushright} \large
            \textbf{Cadre du Projet :}\\
            Projet Majeur DevSecOps \& Cloud Native\\
            Région AWS : \texttt{us-east-1}
        \end{flushright}
    \end{minipage}

    \vfill
    {\large \textbf{Date :} Août 2026 | \textbf{Version :} 1.0 (Production-Ready)}
\end{titlepage}

\tableofcontents
\newpage

% ==============================================================================
% 1. INTRODUCTION GÉNÉRALE
% ==============================================================================
\chapter{Introduction Générale}
L'adoption croissante des architectures conteneurisées et des orchestrateurs tels que Kubernetes a permis d'accélérer considérablement le cycle de livraison logiciel. Toutefois, cette transformation s'accompagne d'une complexité accrue qui élargit la surface d'attaque. Les menaces ciblant la couche applicative et réseau (telles que Log4Shell, les injections SQL, ou le SSRF) nécessitent une capacité d'intervention ultra-rapide afin de contenir l'impact d'une compromission.

L'objectif de ce projet est de concevoir, déployer et valider une plateforme Cloud sécurisée sur AWS, hébergeant un cluster Kubernetes (kubeadm) et une stack SecOps d'Incident Response (IR) automatisée. Provisionnée via Terraform (Infrastructure as Code), la plateforme intègre une application web e-commerce (Symfony) et un pipeline de détection par Machine Learning capable d'évaluer le risque en temps réel et d'exécuter des playbooks d'atténuation structurés selon le modèle normatif \textbf{SANS PICERL}.

% ==============================================================================
% 2. CONTEXTE ET ÉTAT DE L'ART
% ==============================================================================
\chapter{Contexte et État de l'Art}
La sécurisation des architectures Cloud Native impose de dépasser le modèle de sécurité périmétrique traditionnel pour adopter une approche **DevSecOps** et **Zero Trust**. Dans un environnement conteneurisé, les incidents doivent être gérés selon un cycle méthodique rigoureux. La norme **SANS PICERL** définit 6 phases essentielles :
\begin{itemize}
    \item \textbf{Preparation :} Configuration des règles, seuils et mécanismes de réaction.
    \item \textbf{Identification :} Ingestion, parsing et qualification du risque de l'événement.
    \item \textbf{Containment :} Isolation immédiate de la menace (blocage IP, quarantaine réseau/pod).
    \item \textbf{Eradication :} Neutralisation de la source d'attaque (rate limiting, redirection Honeypot).
    \item \textbf{Recovery :} Restauration du fonctionnement normal des services (rollout, scale-out).
    \item \textbf{Lessons Learned :} Traçabilité complète, captures forensic et post-mortem.
\end{itemize}
L'intégration du Machine Learning au sein de ce pipeline permet d'analyser les signaux faibles et d'attribuer un score de risque composite avant le déclenchement des actions correctives.

% ==============================================================================
% 3. ANALYSE DES BESOINS
% ==============================================================================
\chapter{Analyse des Besoins}

\section{Besoins Fonctionnels (BF)}
\begin{itemize}
    \item \textbf{BF-1 (Ingestion \& Parsing IP) :} Ingestion continue des logs NGINX Ingress et applicatifs Symfony avec extraction précise de l'IP source réelle ($X-Forwarded-For$).
    \item \textbf{BF-2 (Scoring ML Hybride) :} Évaluation asynchrone du risque sur une échelle de $0$ à $100$ combinant détection de signatures OWASP et analyse statistique comportementale.
    \item \textbf{BF-3 (Orchestration IR PICERL) :} Déclenchement automatique des playbooks d'atténuation (`full_advanced_protection`) pour les incidents de score $\ge 76$ ou critiques.
    \item \textbf{BF-4 (Mesures d'Atténuation) :} Gestion du blocage IP progressif (Redis), rate limiting adaptatif glissant, quarantaine logique de pod, génération de NetworkPolicies K8s, capture forensic et redirection Honeypot.
    \item \textbf{BF-5 (Supervision \& Action Manuel) :} Interface utilisateur interactive Next.js offrant la visualisation du score de risque, le suivi des phases PICERL et l'exécution synchrone immédiate via le bouton \texttt{EXECUTE NOW}.
\end{itemize}

\section{Besoins Non-Fonctionnels (BNF)}
\begin{itemize}
    \item \textbf{BNF-1 (Performance \& Latence) :} Temps d'exécution du pipeline complet de réaction $< 2$ secondes après ingestion de l'événement.
    \item \textbf{BNF-2 (Découpage \& Asynchronisme) :} Isolation des tâches de réaction via Celery et Redis pour préserver la réactivité des API du backend FastAPI.
    \item \textbf{BNF-3 (Reproductibilité) :} Provisionnement 100\% déclaratif de l'infrastructure AWS via Terraform.
\end{itemize}

% ==============================================================================
% 4. CONCEPTION DE L'ARCHITECTURE AWS ET KUBERNETES
% ==============================================================================
\chapter{Conception de l'Architecture AWS et Kubernetes}

\section{Topology Réseau AWS}
L'infrastructure provisionnée dans la région \texttt{us-east-1} repose sur un VPC dédié (\texttt{10.0.0.0/16}) segmenté en 2 sous-réseaux publics et 2 sous-réseaux privés répartis sur deux zones de disponibilité (\texttt{us-east-1a}, \texttt{us-east-1b}). Un Application Load Balancer (ALB) est positionné dans les subnets publics pour recevoir le trafic Internet et le diriger vers le cluster Kubernetes. Une NAT Gateway équipée d'une Elastic IP assure la sortie sécurisée des nœuds privés.

\begin{table}[H]
\centering
\caption{Spécifications des Instances EC2 du Cluster}
\label{tab:ec2}
\begin{tabularx}{\textwidth}{llllX}
\toprule
\textbf{Instance EC2} & \textbf{Type Instance} & \textbf{Subnet} & \textbf{Accès Réseau} & \textbf{Rôle et Service Exécuté} \\
\midrule
\texttt{k8s-bastion} & \texttt{t3.micro} & Public A & IP Publique & Point d'accès administration SSH \\
\texttt{k8s-master} & \texttt{t3.medium} & Privé A & NAT Gateway & Control Plane K8s (API Server, etcd) \\
\texttt{k8s-worker-1} & \texttt{t3.medium} & Privé A & NAT Gateway & Worker 1 : Symfony, FastAPI, Redis \\
\texttt{k8s-worker-2} & \texttt{t3.medium} & Privé B & NAT Gateway & Worker 2 : Service ML, Celery, Frontend \\
\bottomrule
\end{tabularx}
\end{table}

\section{Architecture du Cluster Kubernetes \& Routage du Trafic}
Toutes les EC2 s'exécutent sous **Ubuntu Server 22.04 LTS**. Le cluster comprend 1 Control Plane et 2 Worker Nodes initialisés via \texttt{kubeadm} avec le runtime \texttt{containerd} (\texttt{SystemdCgroup = true}). Le plugin CNI déployé est **Flannel** (Pod CIDR \texttt{10.244.0.0/16}). L'exposition externe s'effectue via **NGINX Ingress Controller v1.10.0** relié à l'ALB via le NodePort \texttt{32080}.

Le flux du trafic utilisateur suit la séquence :
\[
\text{Internet} \longrightarrow \text{AWS ALB} \longrightarrow \text{NodePort 32080} \longrightarrow \text{NGINX Ingress} \longrightarrow \text{Service K8s} \longrightarrow \text{Pod Applicatif}
\]

\paragraph{Avertissement et Limites Techniques :}
\begin{enumerate}
    \item \textbf{CNI Flannel vs NetworkPolicies :} Flannel assure la connectivité inter-pods mais n'assure pas seul l'enforcement des Kubernetes NetworkPolicies au niveau du noyau Linux. Dans l'architecture actuelle, les manifests YAML des NetworkPolicies sont générés et stockés dans Redis à des fins de préparation ; la migration vers \textbf{Calico} ou \textbf{Cilium} est présentée en chapitre 11 comme perspective.
    \item \textbf{Single Control Plane :} Le cluster comporte un seul Master Node. La mise en œuvre d'un Control Plane multi-masters hautement disponible (HA) constitue également une évolution future.
\end{enumerate}

% ==============================================================================
% 5. IMPLÉMENTATION DE L'INFRASTRUCTURE AVEC TERRAFORM
% ==============================================================================
\chapter{Implémentation de l'Infrastructure avec Terraform}
L'utilisation de Terraform permet d'assurer une gestion globale du State et une reproductibilité parfaite des environnements.

\begin{table}[H]
\centering
\caption{Cartographie des Fichiers Terraform du Projet}
\label{tab:terraform}
\begin{tabularx}{\textwidth}{lX}
\toprule
\textbf{Fichier Terraform} & \textbf{Description et Ressources Provisionnées} \\
\midrule
\texttt{provider.tf} & Spécification du provider AWS et de la région \texttt{us-east-1}. \\
\texttt{networking.tf} & Crée le VPC, Subnets publics/privés, IGW, NAT Gateway, EIP et Tables de routage. \\
\texttt{security\_groups.tf} & Définit les Security Groups (\texttt{bastion\_sg}, \texttt{k8s\_master\_sg}, \texttt{k8s\_worker\_sg}, \texttt{alb\_sg}). \\
\texttt{compute.tf} & Provisionne les 4 instances EC2 avec attachement des clés SSH et Userdata scripts. \\
\texttt{loadbalancer.tf} & Configure l'ALB AWS, le Target Group (NodePort \texttt{32080}) et le Listener HTTP (\texttt{80}). \\
\texttt{flowlogs.tf} & Crée le rôle IAM et active l'exportation VPC Flow Logs vers CloudWatch Logs. \\
\texttt{guardduty.tf} & Active le service de détection automatisée des menaces AWS GuardDuty. \\
\texttt{monitoring.tf} & Configuration des alertes et métriques CloudWatch. \\
\texttt{kubeadm\_token.tf} & Génération dynamique du token d'adhésion sécurisé Kubeadm. \\
\texttt{userdata-master.sh} & Script Shell automatisant l'installation de containerd, Kubeadm, Flannel et Ingress NGINX. \\
\texttt{userdata-worker.sh} & Script Shell automatisant la jonction automatique des workers au Master K8s. \\
\bottomrule
\end{tabularx}
\end{table}

% ==============================================================================
% 6. DÉPLOIEMENT DE L'APPLICATION ET DE LA PLATEFORME SECOPS
% ==============================================================================
\chapter{Déploiement de l'Application et de la Plateforme SecOps}

\section{Application E-Commerce Symfony}
L'application web cliente est développée sous PHP/Symfony. Elle contient les fonctionnalités métier majeures : authentification (\texttt{/login}), catalogue (\texttt{/products}), recherche (\texttt{/search}), panier (\texttt{/cart}) et paiement (\texttt{/checkout}). Elle sert de support aux démonstrations d'attaques (SQLi, XSS, Log4Shell, SSRF, Path Traversal, Brute Force).

\section{Plateforme SecOps Découplée}
\begin{itemize}
    \item \textbf{Frontend Next.js (port 3000) :} Tableau de bord interactif React/TailwindCSS exposant le score de risque global, le flux d'incidents et le bouton \texttt{EXECUTE NOW}.
    \item \textbf{Backend FastAPI (port 8000) :} API centralisant l'ingestion de logs, le parsing d'IP, la persistance PostgreSQL et le dispatch Celery.
    \item \textbf{Worker Celery (\texttt{k8s\_celery}) :} Service asynchrone exécutant les playbooks IR sans bloquer les réponses HTTP du backend.
    \item \textbf{Redis 7 (port 6379) :} Storage in-memory rapide servant de broker Celery et de registre d'état de sécurité.
    \item \textbf{PostgreSQL 15 (port 5432) :} Base relationnelle \texttt{k8s\_platform} contenant les tables \texttt{incidents}, \texttt{incident\_actions}, \texttt{incident\_timers}, \texttt{raw\_logs}, \texttt{risk\_scores}, \texttt{blocked\_ips}, \texttt{circuit\_breaker} et \texttt{notification\_log}.
\end{itemize}

% ==============================================================================
% 7. DÉTECTION ET SCORING DU RISQUE (SERVICE ML)
% ==============================================================================
\chapter{Détection et Scoring du Risque}
Le service ML (\texttt{k8s\_ml}, port 8001) repose sur un **moteur hybride de détection et de scoring** associant la recherche de signatures OWASP et une analyse statistique comportementale.

Le score de risque global ($0 - 100$) est calculé selon la formule composite :
\begin{equation}
\text{Score Global} = 0.50 \times \text{Score\_Sécurité} + 0.30 \times \text{Score\_Fiabilité} + 0.20 \times \text{Score\_Fréquence}
\end{equation}

Les caractéristiques analysées comprennent : le contenu de la requête, les indicateurs malveillants (\texttt{' OR}, \texttt{<script>}, \texttt{jndi:ldap}), la méthode HTTP, le code de statut HTTP, la longueur et l'entropie du log, la fréquence des requêtes par IP source, le namespace et le pod ciblé.

% ==============================================================================
% 8. AUTOMATISATION DE L'INCIDENT RESPONSE (PICERL)
% ==============================================================================
\chapter{Automatisation de l'Incident Response}
Le système orchestre la réponse selon la méthodologie **SANS PICERL**. Le playbook principal \texttt{full\_advanced_protection} applique les mesures d'atténuation suivantes :

\begin{table}[H]
\centering
\caption{Structure du Playbook \texttt{full\_advanced\_protection} (Méthode PICERL)}
\label{tab:picerl}
\begin{tabularx}{\textwidth}{llX}
\toprule
\textbf{Phase PICERL} & \textbf{Action Technique} & \textbf{Description \& Statut d'Implémentation} \\
\midrule
\textbf{Containment} & Blocage IP Progressif & Écriture Redis (\texttt{ir:blocked\_ip}) : 5 min, 30 min, Blacklist 24h. \\
\textbf{Containment} & NetworkPolicy K8s & Génération du YAML restrictif (\texttt{ir:networkpolicy}) conservé dans Redis. \\
\textbf{Containment} & Quarantaine du Pod & Inscription du label \texttt{quarantine=active} et Taint enregistrés dans Redis. \\
\textbf{Eradication} & Rate Limit Adaptatif & Fenêtre glissante Redis (\texttt{ir:ratelimit}) avec backoff exponentiel. \\
\textbf{Eradication} & Redirection Honeypot & Inscription de la règle de dérivation vers Pod piège (\texttt{ir:honeypot}). \\
\textbf{Lessons Learned} & Forensic Capture & Archivage JSON/Redis (\texttt{ir:forensic}) des logs, manifests et pcap (30j). \\
\bottomrule
\end{tabularx}
\end{table}

Playbooks complémentaires disponibles : \texttt{progressive\_block}, \texttt{block\_and\_alert}, \texttt{isolate\_and_block}, \texttt{revoke\_and\_alert}, \texttt{restart\_pod}, \texttt{scale\_out}, \texttt{rollback\_deploy}.

% ==============================================================================
% 9. INTÉGRATION ET FLUX DE DONNÉES
% ==============================================================================
\chapter{Intégration et Flux de Données}
Le pipeline de sécurité s'exécute selon une séquence à 13 étapes :
\begin{enumerate}
    \item NGINX Ingress ou Symfony génère un événement de log.
    \item Le log est émis en JSON à l'API FastAPI via \texttt{POST /logs/ingest}.
    \item FastAPI extrait l'IP source réelle via \texttt{\_extract\_source\_ip} (header \texttt{X-Forwarded-For}).
    \item Le log brut est conservé dans la table \texttt{raw\_logs} de PostgreSQL.
    \item Le service ML évalue l'événement et renvoie le score de risque ($0 - 100$).
    \item Si score $> 40$, un incident est créé ou fusionné dans la table \texttt{incidents}.
    \item Si score $\ge 76$ ou attaque critique, la tâche Celery est déclenchée via \texttt{dispatch\_ir.delay()}.
    \item Le worker Celery exécute le playbook d'atténuation.
    \item Redis enregistre l'état dynamique de sécurité (IPs bloquées, NetworkPolicy, Quarantaine).
    \item PostgreSQL enregistre l'action exécutée dans la table \texttt{incident\_actions}.
    \item Le dashboard Next.js se met à jour en temps réel.
    \item Une alerte est envoyée au Webhook Slack (\texttt{SLACK\_WEBHOOK\_URL}).
    \item Une notification SIEM est transmise via \texttt{SIEM\_WEBHOOK\_URL}.
\end{enumerate}

% ==============================================================================
% 10. VALIDATION EXPÉRIMENTALE ET DÉMONSTRATION
% ==============================================================================
\chapter{Validation Expérimentale et Démonstration}
La validation a été conduite sur l'**Incident \#184** (Attaque Log4Shell lancée depuis l'IP \texttt{32.197.70.124}). 

Lors de l'exécution du playbook \texttt{full\_advanced\_protection}, l'inspection directe de Redis en ligne de commande a confirmé le peuplement de l'ensemble des clés de sécurité :

\begin{lstlisting}[caption={Vérification CLI Redis confirmant l'exécution des mitigations du Playbook}]
$ docker exec k8s_redis redis-cli keys "ir:*"
1) "ir:ratelimit:32.197.70.124:ingress-nginx"
2) "ir:forensic:184-20260520200757"
3) "ir:quarantine:ingress-nginx:ingress-nginx-controller-5458dd5f6-v26dd"
4) "ir:blocked_ip:32.197.70.124"
5) "ir:networkpolicy:ingress-nginx:ingress-nginx-controller-5458dd5f6-v26dd"
6) "ir:honeypot:redirect:32.197.70.124"
\end{lstlisting}

% ==============================================================================
% 11. LIMITES ET PERSPECTIVES
% ==============================================================================
\chapter{Limites et Perspectives}
Afin d'inscrire ce travail dans une démarche d'amélioration continue, les évolutions futures suivantes sont identifiées :
\begin{itemize}
    \item \textbf{Enforcement Réseau complet :} Migration de Flannel vers **Calico** ou **Cilium** (eBPF) pour l'application physique des NetworkPolicies sur le noyau Linux.
    \item \textbf{Détection Runtime eBPF :} Intégration de **Falco** sur les Workers EC2 pour détecter les anomalies au niveau des appels système.
    \item \textbf{Haute Disponibilité K8s :} Passage à un Control Plane multi-masters distribué.
    \item \textbf{Gouvernance \& GitOps :} Déploiement de **Kyverno** (politiques d'admission), **Trivy** (scanning de vulnérabilités) et **ArgoCD** (GitOps).
\end{itemize}

% ==============================================================================
% 12. CONCLUSION GÉNÉRALE
% ==============================================================================
\chapter{Conclusion Générale}
Ce projet d'ingénierie conduit à TEK-UP démontre l'efficacité de l'association entre l'Infrastructure as Code (Terraform), Kubernetes et un pipeline DevSecOps / Incident Response automatisé. Basée sur le cadre normatif SANS PICERL et le Machine Learning, la solution réduit le temps de réaction à moins de 2 secondes tout en assurant une auditabilité complète.

% ==============================================================================
% 13. BIBLIOGRAPHIE
% ==============================================================================
\chapter*{Bibliographie}
\addcontentsline{toc}{chapter}{Bibliographie}
\begin{enumerate}
    \item NIST, \textit{Computer Security Incident Handling Guide}, Special Publication 800-61 Rev. 2, 2012.
    \item Kubernetes Documentation, \textit{Network Policies and Ingress Controllers}, \url{https://kubernetes.io/docs/}, 2024.
    \item HashiCorp, \textit{Terraform AWS Provider Documentation}, \url{https://registry.terraform.io/providers/hashicorp/aws/}, 2024.
    \item SANS Institute, \textit{The Incident Handler's Handbook (PICERL Framework)}, 2019.
\end{enumerate}

% ==============================================================================
% 14. ANNEXES
% ==============================================================================
\chapter*{Annexes}
\addcontentsline{toc}{chapter}{Annexes}
\section*{Annexe A : Exemple de Payload Webhook SIEM}
\begin{lstlisting}[caption={Format JSON transmis au SIEM / SOAR}]
{
  "event_type": "security.incident",
  "timestamp": "2026-05-20T20:07:57Z",
  "incident": {
    "id": 184,
    "type": "log4shell",
    "severity": "CRITICAL",
    "score": 90.3
  },
  "source": { "ip": "32.197.70.124" },
  "actions_taken": ["ip_block", "network_policy", "quarantine", "forensic", "honeypot"]
}
\end{lstlisting}

\end{document}
